\section{Main Theorem}\label{sec:main}

\begin{theorem}[Conditional private direct sum for Cartesian VC-one products]
\label{thm:main}
Under Assumptions~\ref{assump:canonical-product},
\ref{assump:vc-one-factors},
\ref{assump:countably-coded-evaluation}, and
\ref{assump:global-privacy-range}, there are simultaneous universal
constants
\[
C_{\rm up}=65536,
\qquad
C_{\rm quota}=\max\left\{1,K_Y+\frac1{20}\right\},
\qquad
c_{\rm low}>0,
\]
with the following properties.

\begin{enumerate}
\item For every integer
\[
n\ge\left\lceil C_{\rm up}Q_\oplus\right\rceil,
\]
the routed rule \(\bar A_n^{\oplus,Q}\) in
\eqref{eq:prelim-routed-kernel} is a measurable Markov kernel into
\((\mathcal H^\oplus,\mathscr H^\oplus)\), is
replacement-\((\varepsilon,\delta)\)-DP on every pair of adjacent
labeled datasets, and satisfies, for every target \(c\in C\) and every
probability measure \(D\) on \((X,\Sigma)\),
\begin{equation}\label{eq:main-upper-pac}
\Pr_{\substack{S\sim D_c^n\\
                \bar H\sim\bar A_n^{\oplus,Q}(S,\cdot)}}
\left[R_D(h_{\bar H},c)\le\frac1{16}\right]
\ge\frac{15}{16}.
\end{equation}
Consequently,
\begin{equation}\label{eq:main-upper-rate}
\operatorname{SC}_{\varepsilon,\delta}(C)
\le\left\lceil C_{\rm up}Q_\oplus\right\rceil,
\qquad
Q_\oplus
\le C_{\rm quota}
\frac{M_\oplus(C)}{\varepsilon}
\log^2\!\left(
\frac{eM_\oplus(C)}{\varepsilon\delta}
\right).
\end{equation}
This clause holds for every \(0<\delta<1\) allowed by
Assumption~\ref{assump:global-privacy-range}.

\item Fix \(n\ge1\) and suppose
Assumption~\ref{assump:candidate-delta-budget} holds at this candidate.
Every unrestricted measurable replacement-\((\varepsilon,\delta)\)-DP
learner satisfying the universal guarantee \eqref{eq:prelim-sc} at
this \(n\) obeys
\begin{equation}\label{eq:main-candidate-lower}
n\ge c_{\rm low}M_\oplus(C).
\end{equation}
Equivalently, if
\(n<c_{\rm low}M_\oplus(C)\), then for every such learner there are a
full-product target \(c\in C\) and a legal probability measure \(D\)
with \(D(X_i)=\omega_i\) for all \(i\) such that
\begin{equation}\label{eq:main-lower-witness}
\Pr_{\substack{S\sim D_c^n\\
                \Omega\sim A_n(S,\cdot)}}
\left[R_D(h_\Omega,c)>\frac1{16}\right]
>\frac1{16}.
\end{equation}
The learner in this clause may be improper, joint across factors, and
computationally unbounded.

\item Let
\[
n_*=\operatorname{SC}_{\varepsilon,\delta}(C).
\]
If both
\begin{equation}\label{eq:main-nstar-global-cap}
0<\delta\le\frac1{n_*\log(n_*+1)}
\end{equation}
and
\begin{equation}\label{eq:main-nstar-factor-caps}
\delta\le
\frac{c_\delta}
{m_{n_*,i}^2\log(m_{n_*,i}+1)}
\qquad(1\le i\le k)
\end{equation}
hold, then
\begin{equation}\label{eq:main-sandwich}
c_{\rm low}M_\oplus(C)
\le n_*
\le\left\lceil C_{\rm up}Q_\oplus\right\rceil
\le C_{\rm up}C_{\rm quota}
\frac{M_\oplus(C)}{\varepsilon}
\log^2\!\left(
\frac{eM_\oplus(C)}{\varepsilon\delta}
\right).
\end{equation}
If either candidate condition fails, the upper bounds
\eqref{eq:main-upper-rate} remain valid, but no lower bound at \(n_*\)
is asserted.

\item If \(k=1\), then \(M_\oplus(C)=s_1\),
\(Q_\oplus=q_1\), and the upper learner is exactly the unpadded
measurable quotient-first factor rule.  It has zero quotient-risk
residual and satisfies the stronger guarantee
\begin{equation}\label{eq:main-one-factor-upper}
\Pr\left[R_D(h_{\bar H_1},c)\le\frac1{64}\right]
\ge\frac{4095}{4096}.
\end{equation}
For an admissible lower candidate,
\[
\omega_1=1,
\qquad
m_{n,1}=\max\{8,4n\},
\]
the hidden-factor simulation has zero overflow, and the unrestricted
one-factor VC/Littlestone lower conclusion is
\begin{equation}\label{eq:main-one-factor-lower}
n\ge c_{\rm low}\bigl(1+\log^*(d_1+1)\bigr),
\end{equation}
with the strict failure event \eqref{eq:main-lower-witness} on the
contradicted branch.
\end{enumerate}

All utility probabilities are over the iid sample and learner
randomness; privacy is pointwise for adjacent fixed-size datasets and
measurable output events; and every risk is exact distributional binary
\(0\)-\(1\) risk.  The upper horizon is fixed-sample and the lower
horizon is one checked candidate.  The exposed quantities are
\(k,(d_i,s_i,q_i,m_{n,i})_i,M_\oplus(C),Q_\oplus,n,\varepsilon,\delta\).
The fixed quantities are the accuracy and confidence levels
\(1/16,1/64,1/4096\), the equal factor privacy split, the logarithm
convention, and the universal source constants \(K_Y,c_\delta\).
The hidden constants are universal and independent of the class,
factor or quotient cardinalities, distribution supports, block masses,
learner, candidate, and privacy parameters.  No balance, finite-support,
properness, factorwise-output, or computational restriction is hidden.
\end{theorem}

\begin{corollary}[Explicit-rate specialization]
\label{cor:public-rate}
Under Assumptions~\ref{assump:canonical-product},
\ref{assump:vc-one-factors},
\ref{assump:countably-coded-evaluation}, and
\ref{assump:global-privacy-range}, choose the integer quotas
\((q_i)_i\) exactly as in \eqref{eq:prelim-quotas},
\(C_{\rm up}=65536\), and
\(C_{\rm quota}=\max\{1,K_Y+1/20\}\).  Then, for every
\(0<\delta<1\),
\begin{equation}\label{eq:main-public-upper-rate}
\operatorname{SC}_{\varepsilon,\delta}(C)
\le C_{\rm up}C_{\rm quota}
\frac{M_\oplus(C)}{\varepsilon}
\log^2\!\left(
\frac{eM_\oplus(C)}{\varepsilon\delta}
\right).
\end{equation}
If \eqref{eq:main-nstar-global-cap} and
\eqref{eq:main-nstar-factor-caps} hold at the attained sample
complexity \(n_*\), then the left side is also at least
\(c_{\rm low}M_\oplus(C)\).  The probability, horizon, risk, fixed
quantities, and universal hidden-constant dependence are exactly those
declared in Theorem~\ref{thm:main}.
\end{corollary}

\begin{proof}
Proposition~\ref{prop:step-006-public-quota-bridge} verifies the exact
quota choices and absorbs the \(k\) ceiling terms using
\(M_\oplus(C)\ge2k\), \(\varepsilon\le1/10\), and the displayed
logarithmic scale.  Proposition~\ref{prop:step-015-global-upper}
performs the PAC conversion at \(C_{\rm up}Q_\oplus\).
Because \(C_{\rm up}\) and \(Q_\oplus\) are integers,
Proposition~\ref{prop:step-015-conditional-sandwich} proves
\(\lceil C_{\rm up}Q_\oplus\rceil=C_{\rm up}Q_\oplus\), multiplies the
quota bound by \(C_{\rm up}\), and, only after checking both candidate
conditions at \(n_*\), supplies the lower inequality.  These are
precisely the technical checks and probability conversions behind
\eqref{eq:main-public-upper-rate}.
\end{proof}
